\section{Related Work}
\label{submission:sec:related_work}

Our work is positioned at the intersection of model merging, dynamic test-time ensembling, stateful dynamical routing, and the deconstruction of layer-wise depth dynamics in deep neural networks.

\subsection{Static Model Merging}
Model merging has emerged as an active paradigm for combining multiple specialized neural networks into a single multi-task model without retraining. Early methods relied on simple parameter averaging (e.g., in federated learning~\cite{mitchell80, langley00}). Modern static merging techniques perform weight alignment to resolve permutation symmetries, such as GitRe-Basin~\cite{gitrebasin}, or perform task-specific parameter ensembling via Task Arithmetic~\cite{task_arithmetic}. More advanced techniques, such as TIES-Merging~\cite{ties_merging} and DARE~\cite{dare}, identify and eliminate redundant parameters through sign agreement or random pruning prior to merging. While these static approaches successfully create a single multi-task model, they inevitably suffer from inter-task interference and representation collapse when forced to merge highly dissimilar tasks.

\subsection{Dynamic Model Merging \& Subspace Routing}
To circumvent the representational interference of static merging, dynamic model merging (or test-time ensembling) blends task-specific expert adapters on-the-fly based on individual input queries. This is typically achieved using low-rank adapters (LoRAs) plugged into a shared backbone model. 
SABLE~\cite{sable_2024} introduced the concept of activation-space routing, projecting intermediate activations from an early backbone layer onto task-specific PCA subspaces to obtain real-time, scale-free affinity coordinates that dynamically weigh task experts. SPS-ZCA~\cite{pac_zca_2026} proposed Zero-Shot Centroid Alignment to establish lightweight, ultra-low latency routing. 
While stateless dynamic routers achieve high task ensembling accuracy, they suffer from the \emph{routing jitter paradox}~\cite{sable_2024}. Because ensembling weights are computed sample-wise without temporal context, high-frequency coordinate noise causes the active experts to oscillate wildly. This leads to unstable representations, increased inference entropy, and degraded serving stability.

\subsection{Stateful Routing and Kinetics Modeling}
To address the routing jitter paradox, recent frameworks have integrated stateful memory into the routing mechanism. ChemMerge~\cite{chemmerge_2026} modeled routing trajectories as a non-equilibrium chemical kinetics system using continuous-time ordinary differential equations (ODEs) and Arrhenius-like reaction rates. This biochemically inspired stateful smoothing effectively suppresses jitter but introduces significant "inertial drag," causing substantial lag and accuracy degradation immediately following rapid task switches. 
Momentum-Merge~\cite{momentum_merge_2026} deconstructed ChemMerge's complexity, showing that under standard discretization, its biochemical ODEs are mathematically equivalent to simple layer-wise Exponential Moving Average (EMA) smoothing of the routing weights. 
PAC-Kinetics~\cite{pac_kinetics_2026} advanced this by introducing a learning-theoretic stateful router. It proved the input-to-state stability (ISS) and global asymptotic stability (GAS) of continuous-time kinetics systems, and derived a PAC-Bayesian generalization bound for $\beta$-mixing stochastic processes. To suppress temporal lag during rapid switches, PAC-Kinetics introduced an \emph{Adaptive Online Kinetics} mechanism that dynamically scales down state retention when the cosine similarity of consecutive coordinates is low. 

Despite these advancements, existing stateful routers enforce a strictly \emph{spatial homogeneous} architecture: a single, global ensembling weight vector is shared across all network depths.

\subsection{Layer-wise Depth Dynamics in Neural Networks}
A long line of empirical research in deep learning has shown that layers of a deep neural network exhibit highly heterogeneous behaviors. Early layers capture low-level, generic structural features, middle layers form modular semantic concepts, and late layers represent class-specific logit statistics~\cite{DudaHart2nd, MachineLearningI}. This layer-wise differentiation has motivated various decoupled optimization schemes, block-wise training strategies, and layer-specific learning rates. 

While depth-wise representation dynamics are well-studied in static models, their impact on dynamic, test-time stateful model merging has never been explored. In this work, we deconstruct the spatial homogeneity assumption, demonstrating that decoupling the stateful kinetics along the network's depth is both theoretically natural and empirically superior.
